\section*{Erd\H{o}s Problem 1210 --- GPT 6 Astra Ultra partial attempt}


\subsection*{FORMAL RESTATEMENT}
Does an absolute constant $C$ exist such that for every integer $n\geq2$ and every set $A\subseteq\{1,\ldots,n-1\}$ whose distinct elements are pairwise coprime,
\[\sum_{a\in A}\frac1{n-a}\leq\sum_{\substack{p<n\\p\text{ prime}}}\frac1p+C?\]

\subsection*{CONTEXT AND SCOPE}
This replaces the earlier statement-only record with an elementary mathematical attempt. The arguments below establish only their stated partial results; no novelty or full solution is claimed.

\subsection*{ATTACK PLAN AND WORK}
Change variables to $D=\{n-a:a\in A\}\subseteq[1,n-1]$. Pairwise coprimality
has the exact consequence
\[|D\cap(n+p\mathbb Z)|\le1\qquad\text{for every prime }p.\tag{1}\]
Indeed two members of that residue class would correspond to two distinct
elements of $A$ divisible by $p$. Thus the small denominators cannot be chosen
arbitrarily: for each prime, one residue class has occupancy at most one.

For example, if $n$ is even, at most one denominator is even. Bounding its
contribution by $1/2$ and allowing all odd denominators gives
\[\sum_{a\in A}\frac1{n-a}
\le H_{n-1}-\tfrac12H_{\lfloor(n-1)/2\rfloor}+\tfrac12
=\tfrac12\log n+O(1).\tag{2}\]
More generally, for a fixed squarefree $Q\mid n$, at most one denominator is
divisible by each prime of $Q$. After deleting at most $\omega(Q)$ exceptional
denominators, the rest are coprime to $Q$; summing by residue classes yields
\[\sum_{a\in A}\frac1{n-a}
\le\frac{\varphi(Q)}Q\log n+O_Q(1)+\omega(Q).\tag{3}\]
The residue-class estimate in (3) follows by comparing each progression's
harmonic sum with its integral. The dependence of the error on $Q$ cannot
be ignored when allowing $Q$ to grow with $n$.

\subsection*{VERIFICATION AND REMAINING GAP}
Obtain the prime-harmonic scale with an absolute additive constant uniformly
in $n$. The available occupancy constraints only yield the weaker bounds above;
there is no uniform large divisor $Q$ available for arbitrary $n$.

\subsection*{SOURCES}
https://www.erdosproblems.com/1210

\subsection*{FINAL}
\textbf{UNRESOLVED.} The full source problem remains outside the scope of the proved partial results.

\subsection*{COMPLETION ESTIMATE}
$6\%$. This is a conservative subjective estimate toward the whole problem, not a probability that the argument is correct or a claim that a fixed fraction of the problem has been solved.
